% 使用 XeLaTeX 编译。封面必须保留，学校、队号及姓名仅出现在封面。
\documentclass[bwprint]{gmcmthesis}

% === 额外宏包（cls 未内置的）===
\usepackage[framemethod=TikZ]{mdframed}
\usepackage[section]{placeins}
\usepackage{needspace}
\usepackage{subfig}
\usepackage{float}
\usepackage{algorithm2e}

% === TikZ ===
\usepackage{tikz}
\usetikzlibrary{arrows.meta,positioning,shapes.geometric,fit,calc}
\usepackage{longtable}

% === 引用上标格式 ===
\setcitestyle{super,square,comma}

% === 竞赛信息 ===
\title{[论文标题]}
\baominghao{[参赛队号]}
\schoolname{[学校名称]}
\membera{[成员A]}
\memberb{[成员B]}
\memberc{[成员C]}

\begin{document}

\maketitle

% === 摘要 ===
\begin{abstract}

% 写作提示：摘要最后写；按子问题说明所用方法/模型、关键处理、实际数值与单位、结论和验证，突出真正创新。
% 不复述题面、不堆砌模型名、不放复杂公式和表格；2026 年摘要不超过两页，无需英文摘要。
[中文摘要内容：问题概述 + 每个子问题的方法和数值结果 + 结论]

\keywords{[关键词1]\quad [关键词2]\quad [关键词3]}
\end{abstract}

\pagestyle{plain}

% === 正文：摘要结束后自动换页，不插入目录 ===
% 若使用 AI 辅助分析，应在结果前后注明工具名称、版本/型号、开发机构和版本发布日期。
% 模型、公式应有推导或可核验的引用；不得把无法核验的 AI 输出作为既定事实。
\section{问题重述}

% 写作提示：用自己的话交代真实背景、已知条件与数据、待求目标、约束和输出；对含糊概念作必要澄清。
% 不要整段复制题面，也不要在问题重述中提前给出模型公式和求解结果。下方内容仅为结构示例，须按赛题改写。

\subsection{问题背景}
数学建模是解决实际问题的重要工具。本文研究的问题涉及多维度数据分析与优化决策。

\subsection{问题内容}
\textbf{问题一：}对给定数据集进行分析和预测，要求模型具有泛化能力。

\textbf{问题二：}考虑多因素影响，改进和优化模型，提高预测精度。

\textbf{问题三：}结合实际约束条件，设计最优策略方案。

% 技术路线图、研究框架图和算法流程图一律用 paper-diagram Skill 绘制——
% 优先套用其内置模板，产出可编辑 .drawio + PNG，
% 并用 check_layout.py / export_figure.py 校验后插入正文。

\section{问题分析}

% 写作提示：逐问说明输入、输出、约束、数据特点、关键难点、题型判断、拟采用方法及适用理由，
% 并交代各问之间的依赖关系。本节重点回答“为什么这样做”，通常不展开公式或提前给出结果。

\subsection{问题一的分析}
问题一的关键在于特征工程和模型选择，数据中同时存在线性和非线性特征。

\subsection{问题二的分析}
问题二涉及多变量分析和模型融合技术，需要识别主要影响因素。

\subsection{问题三的分析}
问题三是约束优化问题，需要建立目标函数和约束条件。

\section{模型假设}

% 写作提示：只保留模型确实需要且后文会使用的假设；每条说明依据、简化作用、适用范围和可能影响。
% 下方假设仅为示例，不得无依据地照搬“数据完全可靠、因素独立、系统稳定、变量连续可微”等表述。

为简化问题，本文做出以下基本假设：
\begin{enumerate}
    \item 假设题目所给数据真实可靠，不存在系统性的测量误差；
    \item 假设各因素之间的交互作用可以忽略不计；
    \item 假设系统在短期内处于稳定状态；
    \item 假设所有变量连续可微；
    \item 假设缺失值为随机缺失，可采用插补方法处理。
\end{enumerate}

\section{符号说明}

% 写作提示：只列正文实际使用且反复出现的符号，给出含义、单位、取值范围或下标定义；
% 符号首次出现在正文或公式中时仍需就地解释，并保持全文一致。下表内容仅为排版示例。

\begin{table}[htbp]
\centering
\caption{主要符号说明}
\begin{tabular}{c|c|c}
\toprule
\textbf{符号} & \textbf{含义} & \textbf{单位} \\
\midrule
$N$ & 样本数量 & 个 \\
$M$ & 特征维度 & 维 \\
$X$ & 特征矩阵 & --- \\
$y$ & 目标变量 & --- \\
$R^2$ & 决定系数 & --- \\
\bottomrule
\end{tabular}
\end{table}

\section{问题一的模型建立与求解}
% 写作提示：围绕问题目标依次交代数据与输入输出、模型准备、变量和参数、核心关系式或目标与约束、
% 方法选择依据、算法步骤与参数依据、计算环境、关键结果及其解释。下方模型和数值仅为示例，须按实际计算替换。
\subsection{数据预处理}
采用Z-score标准化处理：$x' = \frac{x - \mu}{\sigma}$。

\subsection{模型建立}
采用多元线性回归作为基础模型：
\begin{equation}
y = \theta_0 + \theta_1 x_1 + \theta_2 x_2 + \cdots + \theta_n x_n + \varepsilon
\end{equation}
同时引入随机森林模型进行对比分析，通过集成多棵决策树提高精度和鲁棒性。

\subsection{求解结果}
随机森林的测试集$R^2$达到0.896，优于线性回归的0.834。

\section{问题二的模型建立与求解}
% 写作提示：说明本问是否复用前问的数据、参数或中间结果，以及新增目标和约束；模型选择必须针对题意。
% 给出可复现的求解步骤，并用关键数值、单位和图表直接回答问题二。下方内容仅为结构与排版示例。
\subsection{特征选择}
采用LASSO回归筛选关键特征：
\begin{equation}
\min_{\theta} \frac{1}{2N} \sum_{i=1}^{N} (y_i - \theta^T x_i)^2 + \lambda \|\theta\|_1
\end{equation}

\subsection{模型优化}
基于筛选的关键特征建立优化预测模型，引入交叉项捕捉交互效应。
优化后测试集$R^2$从0.896提升至0.923，RMSE降低约15\%。

\section{问题三的模型建立与求解}
% 写作提示：明确本问与前文的继承关系，写清新增决策变量、目标、约束或评价标准；
% 说明算法流程、参数依据、现实可行性和适用边界。下方优化模型、算法参数和数值结果仅为示例。
\subsection{优化模型}
将问题建模为约束优化问题：
\begin{align}
\min \quad & f(\mathbf{x}) = \sum_{i=1}^{n} c_i(x_i) \\
\text{s.t.} \quad & g_j(\mathbf{x}) \leq 0, \quad h_k(\mathbf{x}) = 0
\end{align}

\subsection{算法求解}
采用遗传算法（种群规模200，迭代100代）进行求解，算法快速收敛到稳定解。

\subsection{结果}
得到最优策略方案：$x_1=0.75$，$x_2=1.20$，$x_3=0.45$，$x_4=0.88$。

\section{敏感性分析}
% 写作提示：只扰动真正不确定且会影响结论的关键参数，并说明基准值、扰动范围依据、输出变化和失效边界。
% 验证方式应适配题型，可结合留出/外部数据、基线、特殊情形、误差或稳定性；下方比例和结果仅为示例。
\subsection{参数敏感性}
对关键参数在基准值附近$\pm 20\%$范围内进行扰动测试。
各参数敏感性均较低（$R^2$最大变化仅1.5\%），模型具有较强的鲁棒性。

\section{模型评价与推广}
% 写作提示：先按题型选择合适的检验方式，如真实/留出数据、基线或特殊情形、误差、稳定性、敏感性等；
% 优点必须有推导、对照或检验支撑，缺点应落到数据、假设、适用范围或算法能力，改进应逐项对应缺点。
% 不要对所有题型机械套用 $R^2$/RMSE，也不要照搬下方通用评价和推广场景。
\subsection{模型优点}
\begin{enumerate}
    \item 理论基础扎实，推导严谨；
    \item 多模型对比确保最优性；
    \item 可解释性强，计算效率高。
\end{enumerate}

\subsection{模型缺点}
\begin{enumerate}
    \item 对数据质量有一定要求；
    \item 未考虑时间动态性。
\end{enumerate}

\subsection{推广}
可推广至金融风控、工业生产、智能交通和能源管理等领域。


% === 参考文献 ===
% 参考文献必须真实：一律用 paper-search Skill 先 search 检索，再 verify 按 DOI 核对元数据。
% 仅将核验通过的文献整理为下方 \bibitem；禁止凭记忆编造题名、作者、期刊、卷期页。
% paper-search 未启用或网络不可用时，只保留能人工核验的真实文献，宁缺毋滥。
% 按正文首次引用顺序编号；书籍注明引用页码，网络资源注明网址和访问日期，引用程序注明来源。
\begin{thebibliography}{99}
% \bibitem[1]{ref1} 作者. 标题[J]. 期刊, 年份.
\end{thebibliography}

% === 附录 ===
\begin{appendices}
\section{核心代码}
% 写作提示：附录应放与正文一致、完整可运行且能复现关键结果的真实代码，并说明环境、依赖、输入和输出。
% 下方代码仅为 listings 排版示例；正式论文应替换为实际程序，较长代码可用 \lstinputlisting 引入文件。
% 如实际使用 AI 辅助编程，须在程序文件开头保留并如实填写来源注释；未使用则删除示例注释。
\begin{Python}{核心求解程序}
# 本程序及代码是在人工智能工具辅助下完成的。
# 工具名称：[实际使用的工具]；版本/型号：[版本]
# 开发机构/公司：[机构]；版本发布日期：[YYYY-MM-DD]
import numpy as np
import pandas as pd
from sklearn.ensemble import RandomForestRegressor
from sklearn.model_selection import train_test_split
from sklearn.metrics import r2_score

data = pd.read_csv('data.csv')
X = data.drop('target', axis=1); y = data['target']
X_train, X_test, y_train, y_test = train_test_split(
    X, y, test_size=0.2, random_state=42
)
model = RandomForestRegressor(n_estimators=100, random_state=42)
model.fit(X_train, y_train)
print(f'R2: {r2_score(y_test, model.predict(X_test)):.4f}')
\end{Python}

\end{appendices}

\end{document}
